\documentclass[10pt,a4paper]{amsart}
\usepackage[margin=19mm]{geometry}
\usepackage{amsmath,amssymb,mathtools}
\usepackage[scr=rsfs,cal=euler]{mathalfa}
\usepackage{xfrac}
\usepackage[unicode,hidelinks,bookmarks=false]{hyperref}
\newcommand{\ie}{i.e.\ }
\newcommand{\cf}{cf.\ }
\newcommand{\resp}{respectively\ }
\newcommand{\Subsection}{\subsection}
\providecommand{\nobreakdash}{\nobreak\hbox{-}}
\setlength{\emergencystretch}{2em}
\multlinegap0pt
\usepackage{amssymb}
\usepackage{esint}
\usepackage{xcolor}
\usepackage{tikz}
\newtheorem{lemma}{Lemma}
\DeclareMathOperator{\BV}{BV}
\title{An alternative approach to the discrete Shnirelman's inequality}
\author{Martina Zizza}
\date{Source: arXiv:2407.09377v1 (CC BY 4.0)}
\begin{document}
\maketitle

\noindent\emph{Source and licence.} An alternative approach to the discrete Shnirelman's inequality. Version arXiv:2407.09377v1 (CC BY 4.0), distributed by arXiv under the Creative Commons Attribution 4.0 International licence, http://creativecommons.org/licenses/by/4.0/, as recorded in the arXiv licence metadata for this exact version and shown on the abstract page https://arxiv.org/abs/2407.09377v1. Full licence notice: \texttt{licenses/arxiv-2407.09377-CC-BY-4.0.txt}.\par\smallskip

\noindent\textit{Editorial scope.} Counted unit: the article's lemma lem:first:exchange (source0.tex lines 1449--1461). The article's complete original proof of it is delivered with it (source0.tex lines 1462--1591, one \texttt{proof} environment of 130 lines). No mathematics is written, repaired or re-derived: the statement and the proof are the article's own lines, in the article's own notation, and no retained line is substituted or re-flowed.  Retained uncounted as the source-local context the proof needs: lines 531--536.  One declared adaptation: exactly 2 ancillary strings inside the retained spans are omitted (image-inclusion commands and, where the source repeats a label, the repeated label definitions) and no other line differs from the source; the figure environments, with their placement, captions and labels, are kept, and the package records the omitted strings in fidelity receipts bound to the affected bodies.  No retained line contains a citation, so the wrapper carries no bibliography and no bibliography entry.  The retained lines contain the article's own diagram code, which the wrapper typesets with the standard tikz packages; no image file is delivered and the package declares no asset dependency.  Editorial formatting only, in addition: an AMS \texttt{amsart} preamble that restates the article's own declarations and macros used by the retained lines, namely the environments \texttt{lemma} on separate counters, together with the standard packages those lines need.  The retained environments print in this document as Lemma 1 in order of appearance, whereas the article prints the same items with the numbering of its own full document.  The complete native source of the article as distributed in the arXiv e-print for this exact version is bundled unchanged as \texttt{sources/162240/source0.tex}.
\begin{equation}\label{eq:rot:vf:rectangle}
\mathbf v_R(x) = \nabla^\perp V_{a,b} = \begin{cases}
\big( 0, \frac{2b}{a} \big( x_1 - \frac{a}{2} \big) \big) & |x_1| \geq \frac{b}{a} |x_2|, 0 \leq x_1 \leq a, \\
\big( -\frac{2a}{b} \big( x_2 - \frac{b}{2} \big),0 \big) & |x_1| < \frac{b}{a} |x_2|, 0 \leq x_2 \leq b.
\end{cases}
\end{equation}

\begin{lemma}\label{lem:first:exchange}
   Let us fix an array $A\subset\mathcal{R}_N$ of length $\ell(A)=M$ with $M\geq 1$. If we denote $A=[\kappa_1,\dots,\kappa_M]$, then there exists a divergence-free vector field $\mathbf v\in L^1([0,1],L^2([0,1]^\nu))\cap  L^1([0,1],\BV([0,1]^\nu))$ with the following property: if $X:[0,1]^\nu\rightarrow [0,1]^\nu$ denotes the time-1 map of the vector field $\mathbf v$, then there exists a positive constant $C>0$, independent of $M,N$ such that 
    \begin{enumerate}
        \item $\|\mathbf v\|_{L^1_tL^2_x}\leq C \sqrt{N^{-\nu-2}M^2}$;
        \item  the map $X$ translates $\kappa_1$ into $\kappa_M$ (and viceversa), namely \begin{equation*}
        X(x)=\begin{cases} (x_1+MN^{-1}-N^{-1},\dots,x_{\nu-1}, x_\nu ), \quad x\in \text{int}(\kappa_1), \\
        (x_1+N^{-1}-MN^{-1},\dots,x_{\nu-1}, x_\nu) \quad x\in \text{int}(\kappa_M),\\
        x \quad \text{otherwise}.
             \end{cases}
    \end{equation*}
    \end{enumerate} 
   
\end{lemma}

\begin{proof}
    We give the proof for $\nu=2$, the extension to higher dimension being straightforward.  Fix $\epsilon>0$ to be chosen later in order to impose the incompressibility condition on the vector field $\mathbf v$. Call
    \begin{equation*}
        a=\frac{\epsilon}{\epsilon+N^{-1}},\qquad b=aMN^{-1},\qquad c=N^{-1}-b;
    \end{equation*}
    then define the following four sets:
     $$A_\epsilon=\left\lbrace (x,y)\in\mathbb R^2: 0\leq y\leq \epsilon,\quad \frac{y}{a}\leq x\leq \frac{y-b}{-a}\right\rbrace,$$
    $$B_\epsilon=\left\lbrace (x,y)\in\mathbb R^2: MN^{-1}-N^{-1}-\epsilon\leq x\leq MN^{-1}, \quad -ax+b\leq y\leq ax+N^{-1}-b\right\rbrace,$$
   $$C_\epsilon=\left\lbrace (x,y)\in\mathbb R^2: N^{-1}-\epsilon\leq y\leq N^{-1}, \quad -\frac{y-N^{-1}}{a}\leq x\leq \frac{y-c}{a}\right\rbrace,$$
    $$D_\epsilon=\left\lbrace (x,y)\in\mathbb R^2: 0\leq  x\leq N^{-1}+\epsilon, \quad ax\leq y\leq -ax+N^{-1}\right\rbrace.$$
    In each one of these sets we define a shear vector field: 
    \begin{equation}\label{eq:vect:field}
        {\mathbf v_1}(x,y)=\begin{cases}
            \left(-\frac{2}{a}y+\frac{b}{a},0\right),\quad (x,y)\in \text{int}(A_\epsilon), \\
            \\
\left(0,2ax+N^{-1}-2b\right),\quad (x,y)\in \text{int}(B_\epsilon), \\
             \\
             \left(-2\frac{y}{a}+\frac{N^{-1}}{a}+\frac{c}{a},0\right),\quad (x,y)\in \text{int}(C_\epsilon), \\
             \\
             \left(0, 2ax-N^{-1}\right),\quad (x,y)\in \text{int}(D_\epsilon). 
        \end{cases}
    \end{equation}
    It can be proved that, with the choice of
   \begin{equation*}
       \epsilon=\frac{N^{-1}}{M-1},
   \end{equation*}
   the vector field is incompressible.
   %Indeed, take a test function $\varphi\in\mathcal{C}^\infty_c(A)$, with $\text{supp}(\varphi)\subset A_\epsilon\cup B_\epsilon$ (we can limit ourselves to the study of this case by the symmetry of the problem). Then we have
   %\begin{align*}
   %0&=\int_{A_\epsilon\cup B_\epsilon} \mathbf { v_1}\cdot\nabla \varphi dx\\&=
   %\int_{0}^\epsilon\int_{\frac{y}{a}}^{-\frac{y}{a}+\frac{b}{a}}\left(-\frac{2}{a}y+\frac{b}{a}\right)\partial_1\varphi dxdy +\int_{MN^{-1}-N^{-1}-\epsilon}^{MN^{-1}}\int_{-ax+b}^{ax+N^{-1}-b}(2ax+N^{-1}-2b)\partial_2\varphi dydx\\&=\int_{0}^\epsilon\left(-\frac{2}{a}y+\frac{b}{a}\right)\varphi\left(-\frac{y}{a}+\frac{b}{a},y\right) dy -\int_{MN^{-1}-N^{-1}-\epsilon}^{MN^{-1}}(2ax+N^{-1}-2b)\varphi\left(x,-ax+b\right)dx\\&=\int_{MN^{-1}-N^{-1}-\epsilon}^{MN^{-1}}a\left(2x-\frac{b}{a}\right)\varphi\left(x,-ax+b\right) dx -\int_{MN^{-1}-N^{-1}-\epsilon}^{MN^{-1}}(2ax+N^{-1}-2b)\varphi\left(x,-ax+b\right)dx,
   %\end{align*}
   %in particular, since the previous equality holds for any choice of the function $\varphi$, we have to require that
   %\begin{equation*}
      % -b-N^{-1}+2b=0 \longrightarrow \epsilon=\frac{N^{-1}}{M-1}.
   %\end{equation*}
   In particular, we remark that $\epsilon$ decreases as the distance between the two cubes increases. 

   \textbf{Step 1.} Since we need to construct the flow that swaps the first and the last cube of the array, we consider $F=A_\epsilon\cup B_\epsilon\cup C_\epsilon\cup D_\epsilon\setminus F'$, where $F'$ is the blue set of Figure \ref{fig:lemmaS2}.
   \begin{figure}
       \centering
       
       \caption{The flow is performed inside the set $A_\epsilon\cup B_\epsilon\cup C_\epsilon\cup D_\epsilon \setminus F',$ where $F'$ denotes the blue frame.}
       \label{fig:lemmaS2}
   \end{figure}
    We consider the flow associated to the vector field defined above, restricted to the set $F$, that is the solution to the Cauchy Problem
    \begin{equation*}
        \begin{cases}
            \dot X(t,(x,y))= \mathbf{ v_1}(X(t, (x,y))), \\
            X(0,(x,y))=(x,y),
        \end{cases}
        (x,y)\in F.
    \end{equation*}
   We observe that in the time interval $t\in[0,2]$ the flow map $X_{t=2}$ has rotated counterclockwise the set $F$ of an angle $\pi$, leaving the intermediate region $A\setminus F$ fixed (see Figure \ref{fig:lemmaS3}). We estimate the $L^2$-norm of the vector field as follows
   \begin{equation*}
       \|\mathbf {v_1}\|^2_{L^2_x}\leq \|\mathbf v_1\|_\infty^2|F|\leq (4M^{2}N^{-2})4N^{-2},
   \end{equation*}
   where we have used that $|F|=2|\kappa_1|+2\frac{N^{-1}}{M} ((M-2)N^{-1}).$
   \begin{figure}
       \centering
       
       \caption{In the time interval $t\in[0,2]$ the flow has rotated counterclockwise the set $F$ of an angle $\pi$. In order to recover the original permutation we have to counter-rotate the set $F$.}
       \label{fig:lemmaS3}
   \end{figure}
   
   \textbf{Step 2.} Here we define the counter rotation vector field. We call
   $$A'_\epsilon=A_\epsilon\setminus (F'\cup (\kappa_1\cup\kappa_M))=[N^{-1},MN^{-1}-N^{-1}]\times\left[0,\frac{N^{-1}}{M}\right],$$
   $$C'_\epsilon=C_\epsilon\setminus (F'\cup (\kappa_1\cup\kappa_M))=[N^{-1},MN^{-1}-N^{-1}]\times\left[N^{-1}-\frac{N^{-1}}{M},N^{-1}\right].$$
    We consider the vector field
   \begin{equation*}
       \mathbf{ v_2}= \mathbf{ v}_{\kappa_1}+\mathbf{ v}_{\kappa_M}+\mathbf{v}_{A'_\epsilon}+\mathbf{v}_{C'_\epsilon},
   \end{equation*}
   where we have used the rotation vector field (Equation \eqref{eq:rot:vf:rectangle}) in the interior of each one of the four sets $\kappa_1,\kappa_M,A'_\epsilon,C'_\epsilon.$ 
 This vector field rotates, in a interval of time $\Delta t=2$, the four sets counterclockwise of an angle $\pi$, as illustrated in Figure. If we estimate the $L^2$-norm of this vector field we find that
 \begin{equation*}
     \|\mathbf{v_2}\|^2_{L^2_x}\leq 2\|\mathbf v_{\kappa_1}\|^2_\infty N^{-2}+2\|\mathbf v_{A'_\epsilon}\|_\infty^2|A'_\epsilon|\leq 4M^2N^{-4}.
 \end{equation*}
 We collect Step 1 and Step 2 defining the vector field of the statement, that is

 \begin{equation*}
     \mathbf v(t,x,y)=\begin{cases}
         8 \mathbf{v_1}(x,y),\quad t\in\left[0,\frac{1}{4}\right), \\
         8 \mathbf{v_2}(x,y), \quad t\in\left[\frac{1}{4},\frac{1}{2}\right).
     \end{cases}
 \end{equation*}
 We could stop the procedure here, since the flow map $T=X_{t=\frac{1}{2}}$ associated to vector field $\mathbf v$ has the following properties:
 \begin{itemize}
     \item $T(\kappa_1)=\kappa_M$ and $T(\kappa_M)=\kappa_1$ and $T\llcorner_{ \kappa_i}$ with $i=1,M$ is a translation;
     \item $T(\kappa_j)=\kappa_j$ for $j\not=1,M$.
 \end{itemize}
So that the map $T$ is a permutation of the cubes of the array. We point out that $T\llcorner_{\kappa_j}$ with $j\not=1,M$ is not the identity. In steps $3$ and $4$ we show how to overcome this issue. 
     \textbf{Step 3.} In this step and the following we correct the map $T$ of the previous two steps such that
     \begin{equation*}
         T\llcorner_{\kappa_j}=Id,\quad \text{a.e. }x,  \quad j\not=1,M. 
     \end{equation*}
     Accordingly to the reference figure, we define
     \begin{equation*}
         \mathbf{v_3}=\sum_{j=2}^{M-1}  \mathbf v_{\kappa_j}.
     \end{equation*}
     Here the estimate of the $L^2$-norm reads as follows:
     \begin{equation*}
         \|\mathbf{v_3}\|^2_{L^2_x}\leq (M-2)\|v_{\kappa_j}\|^2_\infty |\kappa_j|\leq M^2N^{-4}.
     \end{equation*}
     \textbf{Step 4.}
     For this last step we define
     \begin{equation*}
         \kappa_j^A=\kappa_j\cap A'_{\epsilon},\qquad \kappa_j^C=\kappa_j\cap C'_\epsilon,
     \end{equation*}
     and finally
     \begin{equation*}
         \mathbf{v_4}=\sum_{j=2}^{M-1}\left(\mathbf{v}_{\kappa_j^A}+\mathbf{v}_{\kappa_j^C}+\mathbf{v}_{\kappa_j\setminus(\kappa_j^A\cup\kappa_j^c)}\right).
     \end{equation*}
     Here, the computation of the $L^2$-norm is the same as in Step 3. 

        Finally the vector field of the statement is defined as 
         \begin{equation*}
     \mathbf v(t,x,y)=\begin{cases}
         8 \mathbf{v_1}(x,y), \quad t\in\left[0,\frac{1}{4}\right), \\
         8 \mathbf{v_2}(x,y), \quad 
  t\in\left[\frac{1}{4},\frac{1}{2}\right), \\
         8 \mathbf{v_3}(x,y), \quad t\in\left[\frac{1}{2},\frac{3}{4}\right), \\
         8 \mathbf{v_4}(x,y), \quad t\in\left[\frac{3}{4},1\right).\\
     \end{cases}
 \end{equation*}
        This vector field is clearly $\BV$, divergence-free and satisfies the assumptions of the statement. Moreover, we have that
        \begin{align*}
             \|\mathbf v\|_{L^1_tL^2_x}\leq 2\sum_{h=1}^4\|\mathbf{v_h}\|_{L^2_x},
        \end{align*}
        which, together with the estimates obtained for $\|\mathbf{v_h}\|_{L^2_x}$, gives the statement with constant $C=20$, that clearly depends only on the dimension $\nu=2$. 
\end{proof}

\end{document}
